\documentclass[11pt,a4paper]{article}

\usepackage[margin=18mm]{geometry}
\usepackage{fontspec}
\usepackage{xcolor}
\usepackage[most]{tcolorbox}
\usepackage{microtype}
\usepackage{hyperref}

\setmainfont{DejaVu Sans}
\definecolor{Brand}{HTML}{155E46}
\definecolor{Soft}{HTML}{EDF5F0}
\definecolor{Ink}{HTML}{18211C}
\definecolor{Muted}{HTML}{66736C}

\setlength{\parindent}{0pt}
\setlength{\parskip}{5pt}
\linespread{1.03}
\pagestyle{empty}

\hypersetup{
  pdftitle={Research-Mind One-Page Interview Script},
  pdfauthor={Kaushik Sharma}
}

\begin{document}

\begin{center}
  {\LARGE\color{Ink}\bfseries Research-Mind}\par
  \vspace{1mm}
  {\large\color{Brand}\bfseries One-Page Interview Explanation}\par
  \vspace{1mm}
  {\small\color{Muted}Speak this as it is}
\end{center}

\begin{tcolorbox}[colback=Soft,colframe=Brand,boxrule=.7pt,arc=2mm,left=3mm,right=3mm,top=2mm,bottom=2mm]
Hello, my project is called \textbf{Research-Mind}. It is a multi-stage research application that helps students, researchers, and other users understand any topic using recent and traceable sources.

I started this project from a basic tutorial. The original flow had a Search Agent, a Reader Agent, a Writer Chain, and a Critic Chain. I kept the same simple structure because it is easy to understand and debug, but I improved every stage to make the project useful in a real situation.

The problem I wanted to solve is simple. A search engine gives many links, while a normal chatbot may give a good-looking answer without clear evidence. Research-Mind combines both ideas. It searches live sources, studies the collected information, writes a structured report, and then reviews that report.

The project works in four steps. First, the \textbf{Search Agent} understands the user's question and creates a short search plan. The application then collects recent academic papers from OpenAlex, news from Google News, and optional web results from Tavily.

Second, the \textbf{Reader Agent} reads all the collected abstracts and snippets. It compares the sources and identifies important findings, recent signals, disagreements, and missing evidence. This is better than depending on only one webpage.

Third, the \textbf{Writer Chain} uses the Reader's notes and the original source records to create the final report. Every source receives an ID such as S1 or S2, and the Writer uses these IDs as citations. The prompt also tells the model not to invent sources or statistics and to mention uncertainty when the evidence is weak.

Fourth, the \textbf{Critic Chain} receives the completed report and the same source records. It checks whether important claims are supported, whether citation IDs exist, and whether the report contains bias, missing evidence, contradictions, or overconfident conclusions. It returns a score, strengths, areas to improve, and a final verdict. Gemini is the main model for Search, Reader, and Writer. Groq is the optional independent Critic; if it is not configured, Gemini performs the review. The Critic is still an LLM reviewer, not a guaranteed fact checker, and its feedback is shown separately instead of automatically changing the report.

The frontend is built with Streamlit. The user can select the time range, research depth, language, region, audience, and sources. The application shows progress for all four stages and displays the report, Reader notes, sources, Critic feedback, and Search plan in separate tabs. The user can also download the report and ask follow-up questions.

One important limitation is that the system mainly uses abstracts and snippets. A valid source ID also does not guarantee that the source fully supports every claim. For high-stakes research, the user should open and verify the original sources.

My first future improvement would be a controlled Writer--Critic revision loop: the Critic reviews the draft, the Writer revises it once, and a final validation checks the result. I would also add a Python citation-ID validator, claim-to-source support checking, structured JSON feedback, source-quality scoring, and official full-text APIs. For medical, legal, or financial research, I would keep a human approval step before the final report is used.

Overall, this project shows how I converted a simple tutorial workflow into a practical research product while keeping the code readable, sequential, and easy to explain.
\end{tcolorbox}

\end{document}
